


\section{Modèles dénormalisés}

\begin{frame}
\frametitle{Modèles dénormalisés}

Les bases NoSQL reposent sur des modèles \alert{dénormalisés}.

\vfill

\alert{Notions essentielles}

\begin{redbullet}
  \item \alert{Les données sont auto-décrites.} Le contenu vient avec sa propre
  description.
  
 \item \alert{Structures riches.} Listes,
 enregistrements imbriqués, ensembles.

  \item \alert{Typage flexible.} Les données peuvent être typées (``c'est un entier''),
  et/ou structurées (``cette partie est une liste), et/ou ni l'un ni
  l'autre. 
  
   \item \alert{Sérialisation.} Structure \alert{et} contenu doivent pouvoir être transformés
      ensemble en une chaîne de caractères autonome.
\end{redbullet}

\vfill

Le modèle le plus couramment adopté est JSON.
\end{frame}



\begin{frame}{JSON, Javascript Object Notation}

\begin{redbullet}
\item JavaScript Object Notation
\item Initialement créé pour la sérialisation et l'échange d'objets JavaScript
\item Langage pour l'échange de données structurées 
\item Format texte indépendant du language de programmation utilisé pour le manipuler
\end{redbullet}

\vfill 

Utilisation première : échange de données dans un environnement Web (par exemple applications Ajax)\\

\vfill

En NoSQL, JSON apporte flexibilité,  sérialisation et \alert{autonomie}
de représentation

\vfill
\begin{noteblock}
Avec le modèle dénormalisé, on peut ``pré-calculer'' les jointures. 
On le paye par un risque redondance et la perte de la symétrie d'accès.
\end{noteblock}
\end{frame}

\definecolor{lightgray}{rgb}{.9,.9,.9}
\definecolor{darkgray}{rgb}{.4,.4,.4}
\definecolor{purple}{rgb}{0.65, 0.12, 0.82}


\begin{frame}[fragile]{Les bases de JSON}

Structure de base: paire  \textit{clef-valeur} (\textit{key-value})

\begin{minted}{javascript}
"intitule": "Troubles addictifs"
\end{minted}

\bigskip

\alert{Qu'est-ce qu'une valeur?} On distingue les
 valeurs \alert{atomiques} et les valeurs \alert{complexes} (construites)

\medskip
Valeurs atomiques\,: cha\^ines de caractères (entourées 
par les classiques guillemets anglais (droits)), nombres (entiers, flottants) 
et valeurs booléennes (\texttt{true} ou \texttt{false}).

\begin{minted}{javascript}
"année": 2010
\end{minted}


\begin{minted}{javascript}
"valide": false
\end{minted}


\end{frame}

\begin{frame}[containsverbatim,fragile]{Les bases de JSON\,: objects}

Un \textit{objet} est un ensemble de paires clef-valeur.

\vfill

Au sein d'un ensemble de paires, une clef apparait au plus une fois 
(NB : les types de valeurs peuvent être distincts).

\begin{minted}{javascript}
{"intitule": "Troubles addictifs", "date": "12/01/2021","dept": "72"}
\end{minted}

\vfill

Un objet peut être utilisé comme valeur (dite \textit{complexe}) dans une paire clef-valeur.

\begin{minted}{javascript}
"dept": {
  "code": "2A",
  "nom": "Corse-du-Sud"
}
\end{minted}


\end{frame}

\begin{frame}[containsverbatim,fragile]{Les bases de JSON\,: tableaux}


Un \textit{tableau} (\textit{array}) est une liste 
 de valeurs (dont le type n'est pas forcément le même).

\begin{minted}{javascript}
{"cardio": ["Insufisance cardiaque", 
            "Accident vasculaire", 
            "Embolie"]}
\end{minted}

\vfill

Imbrication sans limite\,: tableaux de tableaux, tableaux d'objets contenant eux-mêmes des tableaux, etc.


\end{frame}

\begin{frame}[containsverbatim,fragile]{Les ``documents'' NoSQL}

Un \textit{document} est un objet contenant toutes les informations sur une entité.
\begin{minted}[baselinestretch=.9,
               fontsize=\small,
               framesep=2mm]{javascript}
{
  "_id": "2",
  "prennom": "aapeli",
  "sexe": "m",
  "neLe": "12/11/1962",
  "pathologies": [
    {
      "intitule": "Cancer bronchopulmonaire",
      "date": "03/02/2021",
      "dept": "75"
      },
      {
      "intitule": "Accident vasculaire",
      "date": "11/10/2023",
      "dept": "75"
      }
      ]
}
\end{minted}

\end{frame}


\begin{frame}[containsverbatim,fragile]{Le risque de la redondance}

Les objets ne sont plus séparés comme en relationnel, mais composés
dans un même document.


\begin{minted}[baselinestretch=.9,
               fontsize=\small,
               framesep=2mm]{javascript}
{
  "_id": "2",
  "prennom": "aapeli",
  "sexe": "m",
  "neLe": "12/11/1962",
  "dept": {
     "code": "2A",
     "nom": "Corse-du-Sud",
     "population": 108765
  }
}
\end{minted}

\vfill
On va répéter la description de la Corse du Sud autant de fois qu'il y a de patients.


\end{frame}


\begin{frame}[containsverbatim,fragile]{Tout n'est pas possible}

On ne peut pas se permettre d'avoir des documents sans limite de taille.
Exemple\,: pour les régions et départements, ça va.

\begin{minted}[baselinestretch=.9,
               fontsize=\small,
               framesep=2mm]{javascript}
{
	"_id": "7",
	"nom": "Corse",
	"depts": [
		{
			"code": "2A",
			"nom": "Corse-du-Sud"
		},
		{
			"code": "2B",
			"nom": "Haute-Corse"
		}
	]
}
\end{minted}

\vfill
Mais si je voulais mettre toutes les occurrences de cas dans cet objet,
il ne serait plus vraiment utilisable.
\end{frame}



\begin{frame}[containsverbatim,fragile]{On a perdu la symétrie}

En relationnel, aucun chemin d'accès aux données n'est privilégié.

\vfill

Avec une modélisation dénormalisée on choisit de privilégier 
certaines entités.

\vfill
\begin{minted}[baselinestretch=.9,
               fontsize=\small,
               framesep=2mm]{javascript}
{
  "_id": "2",
  "prennom": "aapeli",
  "pathologies": [
    {"intitule": "Cancer bronchopulmonaire",
      "date": "03/02/2021", "dept": "75"
    },
    {"intitule": "Accident vasculaire",
      "date": "11/10/2023", "dept": "75"
    }
  ]
}
\end{minted}

\vfill
Facile de construire une analyse centrée patient, très difficile 
de regrouper par pathologie ou département. 
\end{frame}

\begin{frame}[fragile]
\frametitle{Les systèmes utilisant JSON}

Beaucoup de systèmes NoSQL s'appuient sur JSON\,:

\begin{redbullet}
       \item MongoDB
       \item ElasticSearch
       \item Cassandra (avec contraintes supplémentaires)
       \item CouchDb et CouchBase, et beaucoup d'autres
 \end{redbullet}

\vfill

\begin{noteblock}
Il n'y a pas de langage de requêtes normalisé et pas d'interopérabilité.
\end{noteblock}
\end{frame}


\begin{frame}[fragile]
\frametitle{À retenir}

Les données NoSQL  se caractérisent par les aspects suivants.

\begin{redbullet}
  \item Une structure \alert{flexible} et \alert{hiérarchisée}
  \item Une structure dénormalisée, complexe avec imbrication de structures
 \end{redbullet}
 
\vfill
La notion \alert{d'autonomie} (pas d'association entre documents) est essentielle.
\begin{redbullet}
       \item Autonomie? On peut mettre en {\oe}uvre un \alert{passage à l'échelle par
       distribution}. Typique des systèmes NoSQL.
       \item Inconvénient\,: un jour ou l'autre il faudra faire des jointures; 
          les systèmes NoSQL ne sont pas bon pour ça.
 \end{redbullet}

\end{frame}


\begin{frame}{La pratique}

À vous de jouer\,!

\vfill


Rendez-vous sur le site \url{https://deptfod.cnam.fr/bd/tp}, exportez les
deux documents JSON proposés et étudiez-les

\vfill

\alert{Défi}\,: proposez une représentation JSON pour les cas, en 
cherchant à éviter toute nécessité de jointure.

\vfill

Vous pouvez vérifier que votre document est correct avec un validateur JSON
en ligne, comme \url{https://jsonlint.com/}

\end{frame}